\section{Introduction}
\label{sec:intro}

Fine-tuning pre-trained foundation models on downstream tasks has become a cornerstone of modern machine learning, but deploying and serving separate models for every task is computationally prohibitive. This operational bottleneck has catalyzed the field of \emph{model merging} \cite{ilharco2022editing, wortsman23}, which combines multiple task-specific expert models into a single, unified network without the need for joint multi-task retraining.

The predominant merging approach, \emph{Task Arithmetic} \cite{ilharco2022editing}, computes linear parameter differences (task vectors) and averages them in Euclidean space. However, Euclidean averaging lacks geometric justification: it assumes a flat loss landscape and an isotropic parameter space, which often fails in practice, leading to \emph{representation drift} and catastrophic interference \cite{ties2023, saim2026}. To address this, geometric parameterizations like \textbf{Orthogonal Fine-Tuning (OFT)} \cite{oft2023} constrain parameter updates to the orthogonal group $\mathrm{O}(d)$ to preserve representation energy and stabilize features. \textbf{OrthoMerge} \cite{yang2026orthomerge} builds on this by merging orthogonal updates on the associated Riemannian manifold: it solves the Orthogonal Procrustes problem to extract rotational components, projects them to the associated Lie algebra $\mathfrak{so}(d)$ via the inverse Cayley transform, averages them, and maps back. Concurrently, spectral alignment methods like \textbf{Sharpness-Aware Isotropic Merging (SAIM)} \cite{saim2026} show that balancing the singular value spectrum (restoring representational isotropy) reduces multi-task interference and enhances performance in Euclidean spaces.

In this work, we present a rigorous theoretical and empirical investigation into the \textbf{limits of representational isotropy on curved manifolds}. Using the orthogonal group $\mathrm{O}(d)$ and its associated Lie algebra tangent space $\mathfrak{so}(d)$ as a diagnostic environment, we evaluate the feasibility of translating Euclidean-based representational isotropy techniques to geometric manifolds. Specifically, we examine a natural manifold extension—\textbf{RIMO} (\textbf{R}iemannian \textbf{I}sometry-respecting \textbf{M}anifold \textbf{O}perations)—which performs Singular Value Decomposition (SVD) directly on the skew-symmetric generators of $\mathfrak{so}(d)$ to balance rotation magnitudes before Cayley reconstruction.

Our investigation reveals two critical phenomena that dictate the viability and physical boundaries of Riemannian model merging, shifting our focus to a deep theoretical diagnostic study:

\textbf{(1) The Orthogonality Condition:} Manifold merging is highly sensitive to non-orthogonality. In standard networks (non-OFT), Procrustes decoupling yields a high-norm linear residual $\rho$, acting as an unconstrained Euclidean displacement that collapses accuracy ($42.07\%$). Incorporating soft orthogonal regularization ($\lambda_{ortho}=2.0$) keeps parameters near the manifold, minimizing residuals and boosting accuracy to $84.55\%$. Crucially, we show that attempting to bypass orthogonal pre-training via a post-hoc SVD projection of an unconstrained base model onto $\mathrm{O}(d)$ is functionally destructive, collapsing merged performance to $15.00\%$ and proving the necessity of native manifold-respecting models.

\textbf{(2) The Tangent Space Spectral Pitfall:} While Euclidean spectral balancing is linear-safe and robust, Lie-algebraic spectral balancing (RIMO, $t > 1.0$) catastrophically scrambles representations, collapsing accuracy to $13.66\%$. We prove that task-specific updates are low-rank; artificially inflating smaller singular values to force isotropy injects skew-symmetric noise. We formalize this through the \textbf{Kernel Distortion Theorem} and the newly derived \textbf{Spectrum Distortion Theorem}, demonstrating that standard SVD solvers introduce non-symmetric coordinate gauges in multi-dimensional kernels, and that the subsequent projection step required to restore skew-symmetry inevitably distorts the active spectrum. Under the non-linear forward Cayley map, this tangent-space noise is mapped to massive, spurious high-dimensional rotations across inactive planes, destroying model functionality.

These insights establish a clear geometric boundary: operations that are safe and effective in Euclidean space can become highly destructive in tangent spaces of non-linear manifolds due to mapping curvature.

\paragraph{Contributions.} Our primary contributions are:
\textbf{(1)} We present a rigorous study of the limits of representational isotropy on curved manifolds, exposing the \textbf{spectral balancing pitfall in Lie algebras}, and proving that inflating smaller singular values injects destructive rotational noise.
\textbf{(2)} We derive the \textbf{Spectrum Distortion Theorem} to demonstrate the mathematical inconsistency of SVD-based spectral modifications under skew-symmetric projection, and we implement and empirically evaluate real Schur decomposition as a symmetry-preserving, projection-free alternative.
\textbf{(3)} We demonstrate that manifold merging requires native manifold-respecting training, show that post-hoc SVD projection collapses performance, and conduct a live pilot experiment of a hard-constrained orthogonal optimizer (projected Riemannian SGD on the Stiefel manifold) that significantly boosts OrthoMerge to $72.08\%$.
\textbf{(4)} We evaluate SOTA test-time adaptive merging (AdaMerging) under disjoint setups, exposing its catastrophic domain-specific overfitting, and show that our proposed \textbf{rank-preserving spectral pruning} (RIMO-Pruned) recovers robust merged accuracies ($91.49\%$ on MLP and $88.16\%$ on Vision Transformers) without test-time calibration.
